\section{The complete fixed-cofactor middle fibre}

The selected square words in Theorem~\ref{thm:transfer} sit inside a larger
finite fibre which can be computed without guessing a word.

\begin{theorem}[Complete fixed-\(Q\) middle fibre]\label{thm:fixed-Q}
Let
\[
 (p,a,R,u,h,r,s,\kappa,D,\alpha,\beta)
\]
be an original E state, in the notation of the preceding sections, and assume
\(h\equiv3\pmod4\).  Put
\[
 j=\frac{h+1}{4},\qquad
 \mathcal U_{p,h}
 =\{U\in\mathbb Z_{>0}:U\mid j^2,\ h\mid p+4U\}.
\]
For \(U\in\mathcal U_{p,h}\), define
\[
 R_U=\frac{p+4U}{h},\qquad
 a_U=\frac{p+R_U}{4},\qquad
 g_U=\gcd(j,U),
\]
\[
 h_U=\frac{g_U^2}{U},\qquad
 r_U=\frac{U}{g_U},\qquad
 \lambda_U=\frac{j}{g_U},\qquad
 s_U=R_U\lambda_U-r_U.                                 \tag{19}
\]
Then
\[
 \boxed{(a_U,R_U,U,h_U,r_U,s_U,\lambda_U)}
\]
is an original M state at the same prime.  Its middle cofactor is exactly
\[
 Q_U=4h_Ur_U\lambda_U-1=h,
\]
and its ordered reciprocal decomposition is
\[
 \boxed{
 \left(a_U,\;p h_Us_U\lambda_U,\;p h_Ur_U\lambda_U\right).} \tag{20}
\]
The assignment \(U\mapsto(a_U,R_U,U,h_U,r_U,s_U,\lambda_U)\) is a
bijection from \(\mathcal U_{p,h}\) onto all original M states at \(p\)
whose normalized cofactor is \(Q=h\).
\end{theorem}

\begin{proof}
Because \(h\le a<p/2\), \(p>2h\), and \(h\ge3\).  From \(U\mid j^2\),
\[
 4U\le4j^2=\frac{(h+1)^2}{4}<2h(h-1)<p(h-1).
\]
Consequently \(0<R_U<p\).  Since
\(hR_U=p+4U\equiv1\pmod4\) and \(h\equiv3\pmod4\), one has
\(R_U\equiv3\pmod4\).  Thus \(a_U\) is integral and
\(p/4<a_U<p/2\).

Prime by prime, \(U\mid j^2\) implies \(g_U^2/U\in\mathbb Z_{>0}\).
The definitions give
\[
 h_Ur_U=g_U,\qquad h_Ur_U^2=U,\qquad
 h_Ur_U\lambda_U=j.
\]
Moreover,
\[
 R_Uj-U=\frac{R_U(h+1)}4-U
 =\frac{p+R_U}{4}=a_U.
\]
It follows that \(h_Ur_Us_U=a_U\) and
\(r_U+s_U=R_U\lambda_U\).  Since
\(\gcd(r_U,\lambda_U)=1\), a common prime divisor of \(r_U,s_U\) would
divide \(r_U,R_U\), and hence \(hR_U-4U=p\).  It cannot equal \(p\), because
\(0<R_U<p\).  Thus \(\gcd(r_U,s_U)=1\).  Finally
\[
 a_U+U=R_Uj,
\]
so \(R_U\mid a_U+U\), and \(U=h_Ur_U^2\mid a_U^2\).  This proves the
forward construction and \(4h_Ur_U\lambda_U-1=4j-1=h\).

Conversely, take an original M state with cofactor \(Q=h\), normalized as
\[
 a_M=h_Mr_Ms_M,\qquad u_M=h_Mr_M^2,\qquad
 \lambda_M=\frac{r_M+s_M}{R_M}.
\]
Then
\[
 j=\frac{h+1}{4}=h_Mr_M\lambda_M,\qquad
 \frac{j^2}{u_M}=h_M\lambda_M^2\in\mathbb Z_{>0}.
\]
Therefore \(U=u_M\mid j^2\).  Also
\[
 a_M+u_M=h_Mr_M(r_M+s_M)=R_Mj,
\]
whence
\[
 p+4U=4(a_M+u_M)-R_M=R_M(4j-1)=R_Mh.
\]
The primitive normalization gives
\(\gcd(j,U)=h_Mr_M\), and (19) recovers
\(h_M,r_M,\lambda_M,s_M\) exactly.  Distinct \(U\)'s remain distinct as
original middle divisors.  This proves bijectivity.
\end{proof}

\begin{corollary}[Exact negative-square divisor slices]
\label{cor:negative-square-slices}
If an original E state has \(\alpha=-4c^2\) or \(\beta=-4c^2\), then
\[
 h\equiv3\pmod4,\qquad\gcd(h,c)=1,
\]
and, with \(j=(h+1)/4\),
\[
 \alpha=-4c^2:\quad
 \mathcal U_{p,h}=\{U\mid j^2:U\equiv c^2\pmod h\},    \tag{21}
\]
\[
 \beta=-4c^2:\quad
 \mathcal U_{p,h}=\{U\mid j^2:c^2U\equiv j^2\pmod h\}. \tag{22}
\]
The words \(U=c^2\) in alpha and \(U=(j/c)^2\) in beta are available exactly
when \(c\mid j\).  That divisibility is not necessary for the full fixed-\(Q\)
fibre to be nonempty.
\end{corollary}

\begin{proof}
The congruence and coprimality assertions were proved in
Theorem~\ref{thm:transfer}.  In the alpha case,
\(p\equiv-4c^2\pmod h\), so
\(h\mid p+4U\) is equivalent to (21).  In the beta case,
\(-4c^2p\equiv1\pmod h\).  Multiplication of
\(p+4U\equiv0\pmod h\) by \(4c^2\), together with
\(4j\equiv1\pmod h\), gives (22), and every step is reversible because
\(\gcd(4c^2,h)=1\).
\end{proof}

Define the finite residue set
\[
 \Delta_h=\{U\bmod h:U\mid j^2\}
 =\left\{\prod_{\ell^e\parallel j}\ell^{f_\ell}\bmod h:
          0\le f_\ell\le2e\right\}.
\]
The exact emptiness criterion is
\[
 \boxed{\mathcal U_{p,h}=\varnothing
 \Longleftrightarrow -p\,4^{-1}\pmod h\notin\Delta_h.} \tag{23}
\]
Thus the fixed-\(Q\) question is a complete finite divisor-residue fibre, not
merely the selected square test.  Its other \(U\)'s can be nonsquares, so the
square-grade conclusion of Theorem~\ref{thm:transfer} remains specific to its
selected words.

\subsection{Composition with the diagonal and radius-eight returns}

For a hard prime, let \(\mathcal E_p\) be the complete original E-state set,
\[
 \mathcal D_p=\{e\in\mathcal E_p:R=D\},
\]
\[
 \mathcal F_p=\{e\in\mathcal E_p:
 h\equiv3\pmod4,\ \mathcal U_{p,h}\ne\varnothing\}.
\]
After applying the proved diagonal return and every fixed-\(Q\) return, the
exact residual source is
\[
 \boxed{
 \mathcal E_p^{\rm res}
 =\{e\in\mathcal E_p:R\ne D,\ 
 [\,h\not\equiv3\pmod4\ \text{or}\ \mathcal U_{p,h}=\varnothing\,]\}.} \tag{24}
\]
Every state in (24) has
\[
 \boxed{\max(|\alpha|,|\beta|)\ge9.}                   \tag{25}
\]
Indeed, radius at most seven is the diagonal profile
\(\alpha=\beta=-1\).  The sole nonsingular radius-eight profile is
\[
 (\alpha,\beta,h)=(8,-4,11).
\]
Here \(j=3\), and \(U=9\in\mathcal U_{p,11}\); (19) gives
\[
 (h_U,r_U,s_U,\lambda_U)
 =\left(1,3,\frac{p+3}{11},1\right),
\]
\[
 a_U=\frac{3(p+3)}{11},\qquad
 R_U=\frac{p+36}{11}.
\]
Thus the earlier radius-eight map is exactly the beta-square specialization
of Theorem~\ref{thm:fixed-Q}.

If a residual state nevertheless has a negative-square coordinate, it obeys
the stronger finite exclusions
\[
 \alpha=-4c^2\Longrightarrow
 \nexists U\mid j^2\text{ with }U\equiv c^2\pmod h,
\]
\[
 \beta=-4c^2\Longrightarrow
 \nexists U\mid j^2\text{ with }c^2U\equiv j^2\pmod h.
\]
All earlier exact conditions remain: \(h\mid\alpha\beta-1\),
\(p\equiv\alpha\pmod h\), \(p\beta\equiv1\pmod h\),
\((h/p)=-1\), neither coordinate is zero or a positive square, and the
exterior atlas bound is unchanged.

The overlaps are typed exactly.  If one source has both negative-square
coordinates, both descriptions parameterize the same full set
\(\mathcal U_{p,h}\); their selected targets coincide exactly when
\(c_\alpha c_\beta=j\).  If \(\delta\) is the squarefree part of
\((R-1)/2\), the diagonal return coincides with a fixed-\(Q=h\) target exactly
when
\[
 \delta=h,\qquad U=j.
\]
Different \(U\)'s in \(\mathcal U_{p,h}\) always produce different original
middle states.

\subsection{A strict hard-prime improvement}

At
\[
 p=41161,\qquad h=11,\qquad j=3,\qquad c=5,
\]
the original E state
\[
 (a,R,u,h,r,s,\kappa,D,\alpha,\beta)
 =(10318,111,9678284,11,938,1,347829,348767,-100,-348756)
\]
has \(\alpha=-4\cdot5^2\), although \(5\nmid3\).  The state equations are
\[
 a=11\cdot938,\quad u=11\cdot938^2,\quad
 4a-p=111,\quad4u+1=111\cdot348767,
\]
\[
 pr+s=111\cdot347829.
\]
Nevertheless \(U=3\mid9\) and \(p+4U=11\cdot3743\), so
Theorem~\ref{thm:fixed-Q} gives
\[
 (a_U,R_U,U,h_U,r_U,s_U,\lambda_U)
 =(11226,3743,3,3,1,3742,1)
\]
and
\[
 \boxed{\frac4{41161}
 =\frac1{11226}+\frac1{462073386}+\frac1{123483}.}      \tag{26}
\]
Primality is certified by the Lucas witness \(22\):
\[
 41160=2^3\cdot3\cdot5\cdot7^3,\qquad
 22^{41160}\equiv1\pmod{41161},
\]
while the residues at exponents \(41160/q\) for \(q=2,3,5,7\) are
\[
 41160,\quad40429,\quad26424,\quad13707,
\]
and each residue minus one is coprime to \(41161\).

By contrast, at the prime \(p=3049\), the original E state
\[
 (a,R,u,h,r,s,\kappa,D,\alpha,\beta)
 =(806,175,20956,31,26,1,453,479,-144,-448)
\]
has ordered denominators \((806,14043,1113244782)\).  Here \(j=8,c=6\), and
\[
 \{U\bmod31:U\mid64\}=\{1,2,4,8,16\},
\qquad c^2=36\equiv5\pmod{31}.
\]
Hence \(\mathcal U_{3049,31}\) is empty.  Primality is certified by the Lucas
witness \(11\): \(11^{3048}\equiv1\pmod{3049}\), while the residues at
exponents \(3048/q\) for \(q=2,3,127\) are \(3048,2516,3041\), each with
predecessor coprime to \(3049\).  A negative-square coordinate alone does not
force a fixed-\(Q\) return.
